\section{Rebuttal to Liberation Theology}

``Liberation'' is not only the goal of esoterism, but also of the various modern movements. There is a fundamental difference. The understanding of the former is that liberation is an individual task. That is, the esoterist seeks to free himself from his own illusions, attachments, desires, and so on. The goal of modern liberation movements it to define a victim and an oppressor; the victim then seeks to gain freedom from the oppressor.

In his book \textit{Francis of Assisi: A Model for Human Liberation}, the Brazilian liberation theologian \textbf{Leonardo Boff} makes the following point:

\begin{quotex}
The theme of liberation is not new, though it certainly is the strongest impulse in modern culture. Generally, we can state that the history of the past five centuries centers in large part on the process of emancipation.
\end{quotex}


This is an important observation. As the \emph{strongest} impulse, any discourse that may challenge or refute such impulse cannot even be heard. To do so would be to put one on the side of the oppressors, hence the enemy of liberation.

\paragraph{The History of Liberation}


In a brief passage, Fr. Boff lists several of the most prominent historical events of emancipation. The book was published in 1981, so an updated version would probably include some movements that were not even anticipated 30 years ago. This is his list, including its most representative figure:

\begin{itemize}


\item \textbf{Galileo Galilei}. The liberation of \emph{reason} from within the religious totality that impeded the free flight of thought in the discovery of the working mechanisms of the world.

\item \textbf{Rousseau}. The liberation of the \emph{citizen} from the absolutism of the kings, to see the citizen as the real bearer and delegation of political powers.

\item \textbf{Marx}. The liberation of the \emph{proletariat} from capitalist economic domination with the aim of arriving at a socialist society without class distinctions.

\item \textbf{Nietzsche}. The liberation of \emph{life}, shortened and suffocated by the sophistication of metaphysics, morals, and culture.

\item \textbf{Freud}. The liberation of the \emph{psyche} from its interior bonds.

\item \textbf{Marcuse}. The liberation of \emph{industrial man}, reduced to one-dimension by the assembly line productions.

\item \textbf{Feminism}. The liberation of \emph{women} from a patriarchal and male culture, toward a less sexist and more personalistic society.
\end{itemize}


\paragraph{Reform and Revolution}


There are two incompatible ways to evaluate Fr. Boff's liberation movements, given that there is a cosmic order implied by the idea of the Good. The first is reform. When the social structures themselves no longer conform to the idea of the Good, whether due to the corruption or incompetence of the acknowledged authorities, then a reform is necessary. Specifically, this applies to the three spheres of religion, politics, and economics and their respective duties of leading the people to true liberation, peace, and prosperity.

Revolution, on the other hand, is actively opposed to the cosmic order and seeks its overthrow. While reform has a terminus, revolution is never ending, since every vestige of order can never be definitively overturned.

\paragraph{Analysis and Objections}


The history of liberation needs to be analyzed by its own standards. Has it truly led to the intended liberation?

\begin{itemize}


\item \textbf{Reason}. Has reason, liberated from any transcendental roots, provided new and viable answers to the problems of life?

  \begin{itemize}

  
\item It has not solved the problems of the nature and destiny of man
  
\item Rather than liberation from transcendental assumptions, it instead has embraced unprovable and unreasonable assumptions such as materialism, determinism, atheism, and positivism.
  
\item In liberating itself from religious and political authority, it has put itself under the corrupting influence of other authorities, so that so-called objective scientific research serves the interests of those false authorities.
  \end{itemize}


\item \textbf{Citizen}. In overthrowing the absolutism of kings, the citizen then put himself under the arbitrary opinions of the general will. Since most of the citizenry is incompetent to understand complex political, economic, and sociological issues, the citizens become beholden to the influence and propaganda of special interests.

\item \textbf{Proletariat}. The proletariat is hostile to the hierarchical principle. Hence, it has overthrown the rule of throne and altar. However, it has not eliminated inequality. Instead, it promotes performers and athletes as a new hierarchy.

\item \textbf{Life}. While unheeded in his own time, Nietzsche and his ideas are now part of popular culture. Life without a transcendent aim is nothing but the will to power. However, power is finite and its distribution is unequal. We look in vain for the arrival of the superman, and see nothing but the mass man everywhere.

\item \textbf{Psyche}. Freud promised the elimination of neuroses, etc., by the liberation of the ego from the superego and its acceptance of the id. In practice, this means sexual liberation. Instead of freedom, this has led to addiction to pornography and the mass use of sexual imagery to control the population and its purchasing habits. The vast number of people on psychotropic drugs in the Western nations indicates the failure of the liberation of the ego to achieve psychological health.

\item \textbf{Industrial man}. The liberation from work is supposed to lead to multidimensional man, of the sort envisaged by Marx: a farmer by day, a philosopher by night. Instead, the increase in leisure has led to a population satisfied only by vulgar entertainment.

\item \textbf{Women}. While women have been liberated in certain areas, feminism has been part of the atomization of society. The first patriarchy is the father-centered family. This is an organic unit that used to form the basis of society. With atomization, each member of the family is an individual and their relationships to each other are nominal, accepted only as long as the relations are freely chosen by each member. An atomized society is easier to control and leads to mistrust, not to more personalistic relations.
\end{itemize}


\paragraph{True Liberation}


True liberation is never addressed in these historical movements. After historical liberation there is nothing left to do, as was shown by the students in \textit{Kidz, Kulcha, Kreativity}\footnote{\url{https://gornahoor.net/?p=1582}}. Any movement that is opposed to the modern world needs to show how Reform will solve the problems of justice rather than Revolution.


\flrightit{Posted on 2014-05-12 by Cologero}

